\documentclass[
	% Basic information
	StudentName     = 孙敬垚,
	StudentNameEng	= {{Name of Graduate (such as HAN Bao-cheng)}},
	AdvisorName     = {{冯兴杰 教授}},
	AdvisorNameEng	= {{Name of Supervisor and Title (such as Prof. FENG Chang-gen)}},
        AdvisorNameEngSecond = {{Name of Second Supervisor and Title}},
	StudentID		= XXXXX,
	Major           = ,
	ResearchInterest= ,
	Department		= XXXX学院,
	DepartmentEng	= {{Name of Department (or School)}},
	% Time
	SubmitYear		= 202X,
	SubmitMonth		= X,
	SubmitDay		= XX,
	DefenseYear		= 202X,
	DefenseMonth	= X,
	DefenseDay		= XX,
	DefenseDateEng	= {{Date of Submit (such as April 2002)}},
	% Title
	Title           = I2TDM,
	TitleEng        = {{English Title}},
	% Other information
	CLCNumber		= TD235,
	UDCNumber		= 624.152.5
]{cauc_thesis}


% 导入页面命令
\newcommand{\inputpage}[1]{
	\input{settings/#1}
}
\linespread{1.5}


\begin{document}
\inputpage{cover}

% =================摘要==================
\newpage
\setlength{\headsep}{1.224cm}
{\centering\chapter*{\zihao{3}\heiti 摘\quad\quad 要}}
\thispagestyle{fancy}
\setcounter{page}{1}
\pagenumbering{Roman}
\chead{\zihao{-5}中国民航大学硕士学位论文}

\addcontentsline{toc}{chapter}{\texorpdfstring{摘\quad\quad 要}{摘\quad\quad 要}}
\zihao{-4}


\setlength{\baselineskip}{20pt}
% ---------------摘要主体---------------
    文本文本文本文本文本文本文本文本文本文本文本文本文本文本文本文本文本文本文本文本文本文本文本文本文本文本文本文本文本文本文本文本文本文本文本文本文本文本文本文本文本文本文本文本文本文本文本文本文本文本文本文本文本文本文本。
 
    文本文本文本文本文本文本文本文本文本文本文本文本文本文本文本文本文本文本文本文本文本文本文本文本文本文本文本文本文本文本文本文本文本文本文本文本文本文本文本文本。文本文本文本文本文本文本文本文本文本文本文本文本文本文本文本文本文本文本文本文本文本文本文本文本文本文本文本文本文本文本文本文本文本文本文本文本文本文本文本文本。
\newline
\zihao{-4}{\bfseries 关键词：}\zihao{-4} 关键词1；关键词2；关键词3；关键词4
% ---------------摘要主体-------------

% ==============Abstract==============
\newpage
\setlength{\headsep}{1.224cm}
{\centering\chapter*{\zihao{3}\textbf{Abstract}}}
\thispagestyle{fancy}
\addcontentsline{toc}{chapter}{Abstract}
\zihao{-4}
% -----------Abstract main body----------
    Text text text text text text text text text text text text text text text, text text text text text text text, text text text text text text text text text text text text text text. Text text text text text text text text text text text text text text text, text text text text text text text, text text text text text text text text text text text text text text. Text text text text text text text text text text text text text text text, text text text text text text text, text text text text text text text text text text text text text text. Text text text text text text text text text text text text text text text, text text text text text text text, text text text text text text text text text text text text text text. Text text text text text text text text text text text text text text text, text text text text text text text, text text text text text text text text text text text text text text.

    Second paragraph text text text text text text text text text text text text text text, text text text text text text text, text text text text text text text text text text text text text text. Text text text text text text text text text text text text text text text, text text text text text. Second paragraph text text text text text text text text text text text text text text, text text text text text text text, text text text text text text text text text text text text text text. Text text text text text text text text text text text text text text text, text text text text text.
\newline
\zihao{-4}\textbf{Key words: }\zihao{-4}\songti keyword1; keyword2; keyword3; keyword4;
% -----------Abstract main body----------

% ==================目录==================
\clearpage
\newpage
\fancypagestyle{plain}{
    \pagestyle{fancy}
}
\pagestyle{fancy}
\chead{\zihao{-5}中国民航大学硕士学位论文}
\titlecontents{toc}
[0pt] % 左边缩进
{\vspace{-50pt}\centering\Large\bfseries} % 是上方间距
{}
{}
{\vspace{20pt}} % 下方间距
\phantomsection
% \addcontentsline{toc}{chapter}{\texorpdfstring{目\quad 录}{}}
\tableofcontents

\newpage
\setcounter{page}{1}
\pagenumbering{arabic}


\setlength{\headsep}{0.624cm}
\fancypagestyle{plain}{
    \pagestyle{fancy}
}%在章节页也显示页眉
% **********正文开始**************
\pagestyle{fancy}
\setlength{\baselineskip}{23.4pt}
{\centering\chapter{绪论}}
\setlength{\baselineskip}{23.4pt}
\zihao{-4}
\section{标题}
正文文字正文文字正文文字正文文字正文文字正文文字正文文字正文文字正文文正文文字正文文字正文文字正文文字正文文字正文文字。

正文文字正文文字正文文字正文文字正文文字正文文字正文文字正文文字正文文正文文字正文文字正文文字正文文字正文文字正文文字。

正文文字正文文字正文文字正文文字正文文字正文文字正文文字正文文字正文文正文文字正文文字正文文字正文文字正文文字正文文字。

正文文字正文文字正文文字正文文字正文文字正文文字正文文字正文文字正文文正文文字正文文字正文文字正文文字正文文字正文文字\footnote{\zihao{-5}脚注文本脚注文本脚注文本脚注文本}。

正文文字正文文字正文文字正文文字正文文字正文文字正文文字正文文字正文文正文文字正文文字正文文字正文文字正文文字正文文字。

\begin{table}[!ht]
	\setlength{\abovecaptionskip}{0cm}
	\tabcolsep=1cm
	\renewcommand\arraystretch{1.2}
	\renewcommand{\tablename}{表}
	\caption{表名}
	\addtocounter{table}{-1}
	\renewcommand{\tablename}{Tab.}
	\caption{Title name}
	\label{tab:lable}
	\vspace{-10pt}
	\begin{center}
		\begin{tabular}{p{5cm}<{\centering} p{5cm}<{\centering}}
			\toprule[1.5pt]
			{{载  体  类  型}}  & {{标  志  代  码}}\\
			\midrule[1.5pt]
			{{磁带（magnetic tape）}} & MT\\
			{磁盘（disk）}  & DK\\
			{光盘（CD-ROW）}  & CD\\
			{联机网络（online）}  & OL\\
			\bottomrule[1.5pt]
		\end{tabular}
	\end{center}
	\vspace{-5pt}
	{\zihao{5}\linespread{1.0}\setlength{\parindent}{2em}
		
		注：如有需要可对表格进行注释说明，可省略。
		
		数据来源：如需对数据来源进行说明，请参照此格式，可省略。
	}
	\vspace{-10pt}
\end{table}

% 使用自定义的 thesistable 环境
\begin{thesistable}{表名}{titleName}{
	{
		{载  体  类  型}} & {{标  志  代  码}}\\
}
{
	{光盘（CD-ROM）} & CD\\
	{联机网络（online）} & OL\\
}
{
	注：如有需要可对表格进行注释说明，可省略。
	
	数据来源：如需对数据来源进行说明，请参照此格式，可省略。
}
\end{thesistable}
\zihao{-4}

正文文字正文文字正文文字正文文字正文文字正文文字正文文字正文文字正文文字正文文字正文文字正文文字正文文字正文文字正文文字正文文字正文文字正文文字正文文字正文文字正文文字正文文字正文文字正文文字正文文字正文文字正文文字正文文字正文文字正文文字正文文字正文文字正文文字正文文字正文文字正文文字正文文字正文文字正文文字正文文字正文文字正文文字正文文字正文文字正文文字正文文字正文文字正文文字正文文字正文文字正文文字正文文字正文文字正文文字正文文字正文文字。

\subsection{标题}

正文文字正文文字正文文字正文文字正文文字正文文字正文文字正文文字正文文字正文文字正文文字正文文字正文文字正文文字正文文字正文文字正文文字正文文字正文文字正文文字正文文字正文文字正文文字正文文字。
\begin{equation}
	I=\frac{V}{R}
\end{equation}

正文文字正文文字正文文字正文文字正文文字正文文字正文文字正文文字正文文字正文文字正文文字正文文字正文文字正文文字正文文字正文文字正文文字正文文字正文文字正文文字正文文字正文文字正文文字正文文字。

\begin{figure}[h]
	\begin{center}
		\setlength{\abovecaptionskip}{0cm}
		\includegraphics[width=0.6\textwidth]{fig1.png}\label{fig1}
		\vspace{5pt}
		\renewcommand{\figurename}{图}
		\caption{图名}
		\addtocounter{figure}{-1}
		\renewcommand{\figurename}{Fig.}
		\caption{Figure name}
	\end{center}
	\vspace{-8pt}
	{\zihao{5}\linespread{1.0}\setlength{\parindent}{2em}
		
		注：如有需要可对图片进行注释说明，可省略。
		
		图片来源：如需对图片来源进行说明，请参照此格式，可省略。
	}
	\vspace{-10pt}
\end{figure}

\zihao{-4}
正文文字正文文字正文文字正文文字正文文字正文文字正文文字正文文字正文文字正文文字正文文字正文文字正文文字正文文字正文文字正文文字正文文字正文文字正文文字正文文字正文文字正文文字正文文字正文文字。

\subsubsection{标题}
正文文字正文文字正文文字正文文字正文文字正文文字正文文字正文文字正文文字正文文字正文文字正文文字正文文字正文文字正文文字正文文字正文文字正文文字正文文字正文文字正文文字正文文字正文文字正文文字。

\subsubsection{标题}
正文文字正文文字正文文字正文文字正文文字正文文字正文文字正文文字正文文字正文文字正文文字正文文字正文文字正文文字正文文字正文文字正文文字正文文字正文文字正文文字正文文字正文文字正文文字正文文字正文文字正文文字正文文字正文文字正文文字正文文字正文文字正文文字正文文字正文文字正文文字正文文字正文文字。

\begin{figure}[h]
	\begin{center}
		\setlength{\abovecaptionskip}{0cm}
		\subfloat[\zihao{5}分图1]{\includegraphics[width=0.5\textwidth]{fig2a.png}\label{fig2a}}\hspace{30pt}
		\subfloat[\zihao{5}分图2]{\includegraphics[width=0.5\textwidth]{fig2b.png}}
		\renewcommand{\figurename}{图}
		\caption{图名}
		\addtocounter{figure}{-1}
		\renewcommand{\figurename}{Fig.}
		\caption{Figure name}
	\end{center}
	\vspace{-10pt}
	% 图注位置
	\vspace{-10pt}
\end{figure}

正文文字正文文字正文文字正文文字正文文字正文文字正文文字正文文字正文文字正文文字正文文字正文文字正文文字正文文字正文文字正文文字正文文字正文文字正文文字正文文字正文文字正文文字正文文字正文文字。

正文文字正文文字正文文字正文文字正文文字正文文字正文文字正文文字正文文字正文文字正文文字正文文字正文文字正文文字正文文字正文文字正文文字正文文字正文文字正文文字正文文字正文文字正文文字正文文字。文文字正文文字正文文字正文文字正文文字正文文字正文文字正文文字正文文字正文文字正文文字正文文字正文文字正文文字正文文字正文文字正文文字正文文字正文文字正文文字正文文字正文文字正文文字正文文字。

正文文字正文文字正文文字正文文字正文文字正文文字正文文字正文文字正文文字正文文字正文文字正文文字正文文字正文文字正文文字正文文字正文文字正文文字正文文字正文文字正文文字正文文字正文文字正文文字。

正文文字正文文字正文文字正文文字正文文字正文文字正文文字正文文字正文文字正文文字正文文字正文文字正文文字正文文字正文文字正文文字正文文字正文文字正文文字正文文字正文文字正文文字正文文字正文文字。

正文文字正文文字正文文字正文文字正文文字正文文字正文文字正文文字正文文字正文文字正文文字正文文字正文文字正文文字正文文字正文文字正文文字正文文字正文文字正文文字正文文字正文文字正文文字正文文字。

正文文字正文文字正文文字正文文字正文文字正文文字正文文字正文文字正文文字正文文字正文文字正文文字正文文字正文文字正文文字正文文字正文文字正文文字正文文字正文文字正文文字正文文字正文文字正文文字正文文字正文文字正文文字。

\subsection{标题}
正文文字正文文字正文文字正文文字正文文字正文文字正文文字正文文字正文文字正文文字正文文字正文文字正文文字正文文字正文文字正文文字正文文字正文文字正文文字正文文字正文文字正文文字正文文字正文文字正文文字正文文字正文文字正文文字正文文字正文文字正文文字正文文字正文文字正文文字正文文字正文文字正文文字正文文字正文文字正文文字正文文字正文文字正文文字正文文字正文文字正文文字。

\begin{table}[!ht]
	\setlength{\abovecaptionskip}{0cm}
	\tabcolsep=1cm
	\renewcommand\arraystretch{1.2}
	\renewcommand{\tablename}{表}
	\caption{表名}
	\addtocounter{table}{-1}
	\renewcommand{\tablename}{Tab.}
	\caption{Table name}
	\label{tab:lable}
	\begin{center}
		\begin{tabular}{ccccc}
			\toprule[1.5pt]
			编号	&长	&宽	&高	&检测结果\\
			\midrule[1.5pt]
			10001	&104	&98	&103	&合格\\
			10001	&104	&98	&103	&合格\\
			10001	&104	&98	&103	&合格\\
			10001	&104	&98	&103	&合格\\
			10001	&104	&98	&103	&合格\\
			10001	&104	&98	&103	&合格\\
			10001	&104	&98	&103	&合格\\
			\bottomrule[1.5pt]
		\end{tabular}
	\end{center}
\end{table}

\begin{table}[!h]
	\addtocounter{table}{-1}
	\setlength{\abovecaptionskip}{0cm}
	\tabcolsep=1cm
	\renewcommand\arraystretch{1.2}
	\renewcommand{\tablename}{表}
	\caption{表名（续）}
	\addtocounter{table}{-1}
	\renewcommand{\tablename}{Tab.}
	\caption{Table name (continued)}
	\label{tab:lable}
	\begin{center}
		\begin{tabular}{ccccc}
			\toprule[1.5pt]
			编号	&长	&宽	&高	&检测结果\\
			\midrule[1.5pt]
			10001	&104	&98	&103	&合格\\
			10001	&104	&98	&103	&合格\\
			10001	&104	&98	&103	&合格\\
			10001	&104	&98	&103	&合格\\
			10001	&104	&98	&103	&合格\\
			10001	&104	&98	&103	&合格\\
			10001	&104	&98	&103	&合格\\
			10001	&104	&98	&103	&合格\\
			\bottomrule[1.5pt]
		\end{tabular}
	\end{center}
\end{table}
\vspace{-16pt}

\subsection{标题}
正文文字正文文字正文文字正文文字正文文字正文文字正文文字正文文字正文文字正文文字正文文字正文文字正文文字正文文字正文文字正文文字正文文字正文文字正文文字正文文字正文文字正文文字正文文字正文文字正文文字正文文字正文文字正文文字。

\section{标题}
正文文字正文文字正文文字正文文字正文文字正文文字正文文字正文文字正文文字正文文字正文文字正文文字正文文字正文文字正文文字正文文字正文文字正文文字正文文字正文文字正文文字正文文字正文文字正文文字。

\section{标题}
正文文字正文文字正文文字正文文字正文文字正文文字正文文字正文文字正文文字正文文字正文文字正文文字正文文字正文文字正文文字正文文字正文文字正文文字正文文字正文文字正文文字正文文字正文文字正文文字正文文字正文文字正文文字正文文字正文文字正文文字正文文字正文文字正文文字正文文字正文文字正文文字正文文字正文文字正文文字。


{\centering\chapter{标题}}

\section{标题}
正文文字正文文字正文文字正文文字正文文字正文文字正文文字正文文字正文文字正文文字正文文字正文文字正文文字正文文字正文文字正文文字正文文字正文文字正文文字正文文字正文文字正文文字正文文字正文文字。

\subsection{标题}
正文文字正文文字正文文字正文文字正文文字正文文字正文文字正文文字正文文字正文文字正文文字正文文字正文文字正文文字正文文字正文文字正文文字正文文字正文文字正文文字正文文字正文文字正文文字正文文字。

正文文字正文文字正文文字正文文字正文文字正文文字正文文字正文文字正文文字正文文字正文文字正文文字正文文字正文文字正文文字正文文字正文文字正文文字正文文字正文文字正文文字正文文字正文文字正文文字。

正文文字正文文字正文文字正文文字正文文字正文文字正文文字正文文字正文文字正文文字正文文字正文文字正文文字正文文字正文文字正文文字正文文字正文文字正文文字正文文字正文文字正文文字正文文字正文文字。


{\centering\chapter{标题}}
\section{标题}
正文文字正文文字正文文字正文文字正文文字正文文字正文文字正文文字正文文字正文文字正文文字正文文字正文文字正文文字正文文字正文文字正文文字正文文字正文文字正文文字正文文字正文文字正文文字正文文字。

\subsection{标题}
正文文字正文文字正文文字正文文字正文文字正文文字正文文字正文文字正文文字正文文字正文文字正文文字正文文字正文文字正文文字正文文字正文文字正文文字正文文字正文文字正文文字正文文字正文文字正文文字。

正文文字正文文字正文文字正文文字正文文字正文文字正文文字正文文字正文文字正文文字正文文字正文文字正文文字正文文字正文文字正文文字正文文字正文文字正文文字正文文字正文文字正文文字正文文字正文文字。

正文文字正文文字正文文字正文文字正文文字正文文字正文文字正文文字正文文字正文文字正文文字正文文字正文文字正文文字正文文字正文文字正文文字正文文字正文文字正文文字正文文字正文文字正文文字正文文字。

{\centering\chapter{标题}}
\section{标题}
正文文字正文文字正文文字正文文字正文文字正文文字正文文字正文文字正文文字正文文字正文文字正文文字正文文字正文文字正文文字正文文字正文文字正文文字正文文字正文文字正文文字正文文字正文文字正文文字。

\subsection{标题}
正文文字正文文字正文文字正文文字正文文字正文文字正文文字正文文字正文文字正文文字正文文字正文文字正文文字正文文字正文文字正文文字正文文字正文文字正文文字正文文字正文文字正文文字正文文字正文文字。

正文文字正文文字正文文字正文文字正文文字正文文字正文文字正文文字正文文字正文文字正文文字正文文字正文文字正文文字正文文字正文文字正文文字正文文字正文文字正文文字正文文字正文文字正文文字正文文字。

正文文字正文文字正文文字正文文字正文文字正文文字正文文字正文文字正文文字正文文字正文文字正文文字正文文字正文文字正文文字正文文字正文文字正文文字正文文字正文文字正文文字正文文字正文文字正文文字。

{\centering \chapter *{结\quad\quad 论}}
\addcontentsline{toc}{chapter}{\texorpdfstring{结\qquad 论}{结\qquad 论}}
正文文字正文文字正文文字正文文字正文文字正文文字正文文字正文文字正文文字正文文字正文文字正文文字正文文字正文文字正文文字正文文字正文文字正文文字正文文字正文文字正文文字正文文字正文文字正文文字。

正文文字正文文字正文文字正文文字正文文字正文文字正文文字正文文字正文文字正文文字正文文字正文文字正文文字正文文字正文文字正文文字正文文字正文文字正文文字正文文字正文文字正文文字正文文字正文文字。

%	------------致谢-------------
{\centering \chapter *{致\qquad 谢}}

\addcontentsline{toc}{chapter}{\texorpdfstring{致\qquad 谢}{致\qquad 谢}}
正文文字正文文字正文文字正文文字正文文字正文文字正文文字正文文字正文文字正文文字正文文字正文文字正文文字正文文字正文文字正文文字正文文字正文文字正文文字正文文字正文文字正文文字正文文字正文文字。

正文文字正文文字正文文字正文文字正文文字正文文字正文文字正文文字正文文字正文文字正文文字正文文字正文文字正文文字正文文字正文文字正文文字正文文字正文文字正文文字正文文字正文文字正文文字正文文字。


\addcontentsline{toc}{chapter}{参考文献}%将“参考文献加入目录中”
\nocite{*}%如果需要只显示引用的参考文献请删除这行
{\centering \bibliography{reference.bib}}

{\centering \chapter *{附\qquad 录}}
\addcontentsline{toc}{chapter}{\texorpdfstring{附\qquad 录}{附\qquad 录}}


{\centering \chapter*{作者简介}}
\addcontentsline{toc}{chapter}{作者简介}

卞兴鹏，是一位在学术领域崭露头角的青年才俊。他本科阶段在华北理工大学打下了坚实的专业基础，本科期间，他勤奋刻苦，积极参与各类专业课程学习与实践活动，凭借优异的成绩和对专业知识的深入钻研，为后续的学术发展积累了宝贵的经验。

随后，卞兴鹏在中国民航大学开启了硕士阶段的学习生涯。在中国民航大学浓厚的学术氛围中，他不断提升自我，展现出了卓越的学术潜力。

在学术研究方向上，卞兴鹏聚焦于 I2TDM 和 LGTDM 领域，致力于通过深入研究，为相关领域的发展贡献力量。他的研究并非一帆风顺，面对复杂的研究课题和技术难题，他始终保持着坚韧不拔的毅力和勇于探索的精神。在无数个日夜中，他沉浸于实验室和文献资料里，反复进行数据的分析与实验的验证。

2023 年，在首届中国智能交通创新挑战赛中，卞兴鹏作为 “CAUCoding” 战队的队长，带领团队在 “基于 ADS - B 的飞行流量预测算法” 赛题上取得了令人瞩目的成绩，荣获二等奖。在这次比赛中，他充分发挥了自己的领导才能和专业优势，组织团队成员共同攻克技术难关，将自己在 I2TDM 和 LGTDM 方面的研究成果灵活运用到实际问题的解决中。这次获奖不仅是对他个人学术能力的高度认可，也彰显了他在团队协作和实际应用方面的出色能力。

卞兴鹏还是冯兴杰教授的亲传大弟子。冯兴杰教授在数据仓库、数据挖掘与智能信息处理等领域有着深厚的造诣，主持和参与了多项国家级、省部级科研项目，在专业期刊和国际学术会议上发表了大量高质量的论文，并且荣获多项省部级科技奖和教学成果奖。在冯兴杰教授的悉心指导下，卞兴鹏在学术道路上少走了许多弯路。冯兴杰教授严谨的治学态度、丰富的科研经验和前瞻性的学术视野，都深深影响着卞兴鹏，激励他不断追求卓越，在学术研究的道路上稳步前行。

目前，卞兴鹏依然活跃在学术前沿，持续在 I2TDM 和 LGTDM 领域进行深入探索，不断追求学术上的突破，为相关领域的发展注入新的活力。
\end{document}